
\begin{recipe}{German Chocolate Cake}\tag{cake}
  \rating    5
  \nutrition 2
  \health    2
  \workload{medium}
  \time{60}
  \yield{9x13 inch cake}
  \servings{18--24 pieces}
  \source{Dean}
  \maketitle
  % \photo{}

  \begin{ingredients}
    1 & Chocolate \reciperef{Box Cake}\\
    & \reciperef{Coconut Pecan Frosting}
  \end{ingredients}

  Bake in greased $9\times13$\inch baking pan. Allow cake to cool completele. Split
  cake horizontally and frost while frosting is still warm.
\end{recipe}

\begin{recipe}{Coconut Pecan Frosting}\tag{frosting}
  \rating    5
  \nutrition 2
  \health    2
  \workload{medium}
  \time{20}
  \yield{frosting for 9x13 inch cake}
  \source{Dean}
  \maketitle
  % \photo{}

  \begin{ingredients2}
    12 \oz & evaporated milk\\
    5 & egg yolks\\
    1\half \c & sugar\\
    \threequarter \cup & butter\\
    2\half \cup & sweetened coconut\\
    2 \cup & pecans, chopped fine\\
    1 \T & vanilla
  \end{ingredients2}

  Stir together milk, sugar, butter, and eggs in large sauce pan. Stirring
  continuously, bring to temperature and cook at low boil for 15 minutes. Stir in
  rest of ingredients and cool. Frost while still partially warm.

  \begin{note}
    This recipe will give plenty of frosting, likely with a bit left over.
  \end{note}
\end{recipe}


%%% Local Variables:
%%% mode: latex
%%% TeX-master: "../Cookbook"
%%% End:
